\subsection{Limiti unilateri e vincoli di dominio}
\label{sec:esercizi-limiti-limiti-unilateri-e-dominio}

\begin{exercise}[Limiti unilateri con polo semplice]
\label{ex:limiti-foglio-2311-07}
\difficoltaesercizio{0}
Calcolare il limite
\[
  \lim_{x\to1}\frac{x^2-6x+5}{x^2-2x+1}.
\]
\rimandoteoria{def:lez14-limiti-unilateri}
\end{exercise}
\begin{proof}[Soluzione dell'esercizio]
Per $x\ne1$, fattorizzando e semplificando si ottiene
\[
  \frac{x^2-6x+5}{x^2-2x+1}
  =\frac{(x-1)(x-5)}{(x-1)^2}
  =\frac{x-5}{x-1}.
\]
Poiché $x-5\to-4<0$ per $x\to1$ e $x-1$ cambia segno,
\[
  \lim_{x\to1^-}\frac{x-5}{x-1}=+\infty,
  \qquad
  \lim_{x\to1^+}\frac{x-5}{x-1}=-\infty.
\]
I limiti unilateri sono diversi, quindi il limite bilatero non esiste.
\end{proof}

\begin{exercise}[Fattore ripetuto al denominatore]
\label{ex:limiti-foglio-2311-09}
\difficoltaesercizio{0}
Calcolare il limite
\[
  \lim_{x\to4}\frac{x^2-10x+24}{(x-4)^2}.
\]
\rimandoteoria{def:lez14-limiti-unilateri}
\end{exercise}
\begin{proof}[Soluzione dell'esercizio]
Per $x\ne4$, fattorizzando e semplificando si ottiene
\[
  \frac{x^2-10x+24}{(x-4)^2}
  =\frac{(x-4)(x-6)}{(x-4)^2}
  =\frac{x-6}{x-4}.
\]
Poiché $x-6\to-2<0$ per $x\to4$ e $x-4$ cambia segno,
\[
  \lim_{x\to4^-}\frac{x-6}{x-4}=+\infty,
  \qquad
  \lim_{x\to4^+}\frac{x-6}{x-4}=-\infty.
\]
I limiti unilateri sono diversi, quindi il limite bilatero non esiste.
\end{proof}

\begin{exercise}[Sostituzione in un quoziente con radice]
\label{ex:limiti-foglio-2311-10}
\difficoltaesercizio{0}
Calcolare il limite
\[
  \lim_{x\to1}\frac1{\sqrt{x}-2}.
\]
\rimandoteoria{thm:lez16-limite-sostituzione}
\end{exercise}
\begin{proof}[Soluzione dell'esercizio]
Poiché $\sqrt{x}-2\to-1\ne0$ per $x\to1$,
per le proprietà dei limiti si ha
\[
  \lim_{x\to1}\frac1{\sqrt{x}-2}
  =\frac1{1-2}=-1.
\]
\end{proof}

\begin{exercise}[Denominatore quadratico di segno positivo]
\label{ex:limiti-foglio-2311-11}
\difficoltaesercizio{0}
Calcolare il limite
\[
  \lim_{x\to1}\frac2{(x-1)^2}.
\]
\rimandoteoria{def:lez14-limiti-unilateri}
\end{exercise}
\begin{proof}[Soluzione dell'esercizio]
Per $x\ne1$ si ha $(x-1)^2>0$ e $(x-1)^2\to0$ per $x\to1$.
Pertanto
\[
  \lim_{x\to1^-}\frac2{(x-1)^2}
  =\lim_{x\to1^+}\frac2{(x-1)^2}=+\infty.
\]
Il limite bilatero è quindi $+\infty$.
\end{proof}

\begin{exercise}[Segno di un quoziente presso zero]
\label{ex:limiti-foglio-2311-15}
\difficoltaesercizio{0}
Calcolare il limite
\[
  \lim_{x\to0}\frac{2x-5}{x^2}.
\]
\rimandoteoria{def:lez14-limiti-unilateri}
\end{exercise}
\begin{proof}[Soluzione dell'esercizio]
Poiché $2x-5\to-5<0$ e $x^2\to0^+$ per $x\to0$,
si ha
\[
  \lim_{x\to0^-}\frac{2x-5}{x^2}
  =\lim_{x\to0^+}\frac{2x-5}{x^2}=-\infty.
\]
Il limite bilatero è quindi $-\infty$.
\end{proof}

\begin{exercise}[Denominatore non nullo nel punto]
\label{ex:limiti-foglio-2311-22}
\difficoltaesercizio{0}
Calcolare il limite
\[
  \lim_{x\to0}\frac{\sin x}{x-\pi}.
\]
\rimandoteoria{thm:lez14-limite-rapporto}
\end{exercise}
\begin{proof}[Soluzione dell'esercizio]
Poiché $\sin x\to0$ e $x-\pi\to-\pi\ne0$ per $x\to0$,
per le proprietà dei limiti si ha
\[
  \lim_{x\to0}\frac{\sin x}{x-\pi}
  =\frac0{-\pi}=0.
\]
\end{proof}
